%% gorti SoftwareX original software publication candidate.
%% Based on SoftwareX's official 2026 original-software LaTeX template.
\documentclass[preprint,12pt,a4paper]{elsarticle}

\usepackage{amssymb}
\usepackage{hyperref}
\usepackage{booktabs}
\setlength{\parindent}{0pt}

\journal{SoftwareX}

\begin{document}

\begin{frontmatter}

\title{gorti: An open-source IEEE 1516-2010 Run-Time Infrastructure in Go}

\author[aff1]{Changbeom Choi\corref{cor1}}
\ead{me@cbchoi.info}
\cortext[cor1]{Corresponding author}
\address[aff1]{Affiliation and full postal address to be confirmed before submission}

\begin{abstract}
The High Level Architecture coordinates independently developed simulation
components, but complete open implementations of its Run-Time Infrastructure
remain limited. gorti is an IEEE 1516-2010 Run-Time Infrastructure implemented
in Go with federate paths for Go, Python, and the C++ Dynamic Link Compatible
interface. It provides federation, declaration, object, time, ownership, data
distribution, synchronization, save/restore, and management-object services.
Deterministic event logs make runs inspectable and repeatable. A fail-closed
cross-RTI harness compares canonical two-process semantics before reporting
paired performance results. gorti supports simulation research, education,
continuous integration, and transparent single-node distributed simulation.
\end{abstract}

\begin{keyword}
High Level Architecture \sep Run-Time Infrastructure \sep distributed
simulation \sep logical time \sep reproducibility \sep interoperability
\end{keyword}

\end{frontmatter}

\section*{Required Metadata}

\section*{Current code version}

\begin{table}[!h]
\small
\begin{tabular}{|l|p{5.1cm}|p{7.4cm}|}
\hline
\textbf{Nr.} & \textbf{Code metadata description} &
\textbf{Current gorti value} \\
\hline
C1 & Current code version & v0.9.0 submission candidate; immutable tag pending \\
\hline
C2 & Permanent GitHub link to code/repository used for this code version &
\url{https://github.com/cbchoi/gorti/tree/v0.9.0}; activate and verify before submission \\
\hline
C3 & Legal Code License & MIT \\
\hline
C4 & Code versioning system used & git \\
\hline
C5 & Software code languages, tools, and services used &
Go, Python, C++, Protocol Buffers, gRPC, MkDocs, GitHub Actions \\
\hline
C6 & Compilation requirements, operating environments \& dependencies &
Go 1.22 or later; Python 3.11 or later for the Python SDK; C++17 and Conan 2
for the C++ SDK; Linux, macOS, or Windows \\
\hline
C7 & Link to developer documentation/manual &
\url{https://github.com/cbchoi/gorti/tree/v0.9.0/docs}; replace with the
verified public documentation URL before submission \\
\hline
C8 & Support email for questions & \href{mailto:me@cbchoi.info}{me@cbchoi.info} \\
\hline
\end{tabular}
\caption{Code metadata. Tag and documentation links are release gates and must
be live before submission.}
\label{tab:metadata}
\end{table}

\section{Motivation and significance}
\label{sec:motivation}

Distributed simulation combines models that may be developed by different
teams, implemented in different languages, and executed in separate processes.
The IEEE High Level Architecture (HLA) defines the framework, federate
interface, and object-model rules for coordinating those models
\cite{ieee1516,ieee15161,ieee15162}. The Run-Time Infrastructure (RTI) is the
middleware that implements federation membership, data exchange, ownership,
and logical-time coordination. Its behavior is consequently part of the
scientific experiment: an ordering defect, an early time grant, or an
undocumented success boundary can alter a run even when model code is
unchanged.

Existing HLA users can choose mature commercial RTIs, while open research
implementations often cover a narrower service surface or embed federation
logic in one process. Commercial software is appropriate for many deployments,
but licensing and implementation opacity can limit unrestricted continuous
integration, algorithm experiments, and archival inspection. A research RTI
therefore needs more than source availability. It needs executable evidence
that identifies which service calls, callbacks, payloads, and logical times
were observed, and it must not label two runs equivalent merely because their
event counts match.

gorti addresses this need with a standalone server and independent federate
processes. It targets IEEE 1516-2010 and exposes federate paths for Go, Python,
and the IEEE C++ Dynamic Link Compatible surface. Its event log uses stable
ordering and canonical encoding for deterministic workloads. The verification
harness projects detailed implementation logs into common federation
management (FM), declaration management (DM), object management (OM), and time
management (TM) records. Missing callbacks, incorrect payloads, wrong logical
times, duplicate delivery, and provenance disagreement all fail validation
before any performance ratio is calculated.

A typical user starts \texttt{rtid}, creates a federation from versioned FOM
modules, and launches two or more independent federates. The federates declare
their publications and subscriptions, exchange objects and interactions, and
coordinate logical time through requests and asynchronous grants. Event logs,
manifests, and metrics can then be used for regression testing or archived as
experimental evidence.

The software contributes five reusable elements: an open Go RTI service
implementation; Go, Python, and C++ federate paths; deterministic lifecycle and
event evidence; canonical cross-language and cross-RTI scenarios; and a paired
AB/BA comparison protocol that couples performance samples to semantic parity.

\section{Software description}
\label{sec:description}

\subsection{Software architecture}

The \texttt{rtid} process exposes versioned Protocol Buffer contracts over
gRPC. Transport handlers delegate to managers for federation membership,
declarations, object instances, interactions, ownership, regions, logical
time, synchronization, save/restore, and management objects. The managers use
an explicit federation generation so delayed callbacks or state from a
destroyed federation cannot enter a later federation created with the same
name.

Each joined federate owns a callback event stream. Object updates and
interactions can be receive ordered or timestamp ordered. Timestamped
recipients are classified under one authoritative time-evaluator reservation.
Required timestamp-order callbacks are delivered before the grant that makes
their timestamp reachable, and a delivery failure withholds the corresponding
grant. Teardown stops admission, drains or cancels event streams, removes
membership, and closes event-log and save resources.

\begin{figure}[!h]
\centering
\fbox{\begin{minipage}{0.92\linewidth}
\centering
Go federate SDK \quad Python asyncio SDK \quad C++ DLC SDK\\[4pt]
$\Updownarrow$ gRPC service calls and callback event streams $\Updownarrow$\\[4pt]
\textbf{rtid transport handlers}\\[3pt]
$\Downarrow$\\[3pt]
Federation $|$ Declaration $|$ Object $|$ Time $|$ Ownership $|$ DDM\\[3pt]
$\Downarrow$\\[3pt]
Deterministic event log $|$ Save/restore $|$ Admin status $|$ Metrics
\end{minipage}}
\caption{gorti process and service architecture. Federates are independent
processes; service managers share generation-aware federation state.}
\label{fig:architecture}
\end{figure}

The Go SDK resolves FOM names to wire handles and emits typed callbacks without
exposing generated Protocol Buffer types to application code. The Python SDK
provides an \texttt{asyncio} API and a 1516-shaped ambassador adapter. The C++
SDK supplies the IEEE 1516.1-2010 DLC header surface on supported Unix-like
platforms. These access paths share the same server service managers.

\subsection{Software functionalities}

Federation management covers creation, joining, synchronization, resignation,
and destruction. Declaration and object management cover publication,
subscription, name reservation, registration, discovery, attribute reflection,
interaction receipt, and deletion. Time management includes regulation,
constraints, lookahead, timestamp-order delivery, and the TAR, TARA, NER,
NMRA, and FQR request families. The conformance program also exercises
ownership, data distribution, message retraction, save/restore, and
management-object services.

Synchronous interaction sends use a persistent client stream and a
caller-visible server acknowledgement. The call returns only after request
serialization, transport, server validation and processing, durable logging
when configured, and the ACK boundary. This makes success and failure
unambiguous. A separate asynchronous API is future work rather than a change to
the established synchronous semantics.

The server provides a loopback admin endpoint for the \texttt{rti-top}
terminal monitor and a Prometheus endpoint. The supported production topology
for this release is a single RTI process; cluster and failover code is not
claimed as production ready.

\subsection{Quality control and reproducibility}

Go unit and integration tests cover managers, protocol handlers, persistence,
SDK behavior, and failure paths. Python tests cover async transport, the public
ambassador surface, generation fencing, and cross-process behavior. C++
compile-time lockfiles protect DLC API shape, while runtime tests cover
standard exception mapping. Cross-language fixtures verify encodings and live
publish/subscribe behavior.

Concurrency-sensitive Go packages run with the race detector. Repeated tests
exercise atomic mixed-recipient fanout, failed-delivery grant withholding,
exactly-once grant evaluation, stream ACK drain, forced cancellation, and
complete federation teardown. Documentation builds in strict mode, and release
metadata is checked across Python, C++, citation, CodeMeta, and Git tag values.

The fair-comparison workload fixes FOM bytes and SHA-256, seed 1516, two
independent federate processes, choreography, callback handling, file server
logging, and measurement boundaries. A claim-grade session uses five warmup
pairs and twenty measured pairs with ten AB and ten BA orientations. The
analyzer validates full canonical FM/DM/OM/TM records and exact delivery
accounting before reporting medians, p95, p99, paired ratios, confidence
intervals, and order effects.

\section{Illustrative examples}
\label{sec:example}

The \texttt{examples/go-tar-wait} program demonstrates federation-wide logical
time. After starting \texttt{rtid}, the user launches two independent
processes, a waiter and a peer. Both become time regulating and time
constrained. The waiter requests logical time 5 while the peer remains at time
0 with lookahead 2. The peer's state constrains the federation lower bound, so
the waiter receives no immediate grant. After a three-second delay, the peer
also requests time 5. The lower bound then permits both
\texttt{TimeAdvanceGrant(5)} callbacks.

\begin{verbatim}
go build -o bin/rtid ./rti/cmd/rtid
go build -o bin/go-tar-wait ./examples/go-tar-wait
./bin/rtid --listen 127.0.0.1:8442 --log-dir ./eventlogs
./bin/go-tar-wait --role waiter
./bin/go-tar-wait --role peer
\end{verbatim}

The waiter prints that \texttt{TAR(5)} is blocked, then reports a passing grant
after the peer advances. Both processes exit, allowing an automated test to
distinguish success from a hung federation. The sequence illustrates three
important properties: a time request is not a local clock assignment; another
regulating federate can delay progress; and grants arrive asynchronously after
the federation-wide safety condition becomes true.

The cross-RTI verification uses a larger two-process example. Both Pitch pRTI
and gorti receive the same FOM path and generated payloads. The scenario
creates and joins a federation, crosses synchronization barriers, publishes
and subscribes, registers and discovers an object, sends timestamped attribute
updates and interactions, advances time, observes exact grants and callbacks,
then resigns and destroys the federation. A canonical projector emits four
records only after every required observation is present.

\section{Impact}
\label{sec:impact}

gorti makes RTI behavior available for inspection and controlled modification.
Researchers can evaluate time-management, ownership, data-distribution, or
transport strategies while retaining a stable semantic gate. Deterministic
logs make ordering regressions visible and support replay analysis. The
standalone process model also allows application federates to be tested in CI
without embedding RTI state in the model process.

The verification method is an additional research contribution. Cross-system
benchmarks can be misleading when one implementation omits logging, uses fewer
processes, handles callbacks differently, or returns success before equivalent
work is complete. gorti's orchestrator records executable hashes, FOM hash,
workload contract, environment, process order, raw logs, accounting, and
statistics. It rejects a session if provenance or canonical semantics differ.
The same pattern can be reused to evaluate future RTI changes or other
distributed middleware with observable event contracts.

The retained submission-candidate experiment used a 12th generation Intel
Core i7-12700 Windows host and the fixed 5+20 AB/BA protocol. The synchronous
\texttt{sendInteraction} medians were 165.600 microseconds for gorti and
121.275 microseconds for Pitch. The paired median gorti/Pitch ratio was 1.218
with a deterministic paired-bootstrap 95\% interval of 1.120--1.327. gorti was
faster at the completed-delivery boundary and competitive for object updates.
These measurements identify a transport/ACK optimization target; they are not
a general vendor ranking.

Microbenchmarks attributed approximately three microseconds to the complete
registry and durable file-log path, while the persistent gRPC request/ACK
round trip was tens of microseconds. Runtime thread tuning, shared write
buffers, and experimental prepared ACK frames did not produce a stable gain
and were rejected. This negative result is useful because it narrows future
work to an explicit API or transport decision instead of weakening ordering or
delivery guarantees.

The project currently provides no defensible public user-count or download
claim. Its demonstrated impact is therefore technical and methodological:
complete executable service scenarios, independent federate processes,
cross-language interfaces, reproducible semantic evidence, and a transparent
baseline for future simulation research. Adoption metrics and publications
using gorti should be added only when independently verifiable.

Limitations are explicit. The release targets IEEE 1516-2010 only and does not
support HLA 1.3 or IEEE 1516-2000. It has no Java SDK and no formal IVCT
certification. Windows C++ support is incomplete. Pitch comparison is
black-box evidence from public API returns, callbacks, and permissible logs;
it cannot establish source-level identity. Proprietary Pitch binaries and
restricted logs are not redistributed.

Future work will keep the strict synchronous API unchanged and evaluate a
separate bounded asynchronous interaction API with backpressure, ordered
completion handles, and explicit failure reporting. It must be compared only
with an equivalent asynchronous Pitch choreography. Timestamp-order delivery
before grants, atomic recipient reservation, generation fencing, teardown, and
delivery accounting remain acceptance gates.

\section{Conclusions}
\label{sec:conclusions}

gorti provides an open IEEE 1516-2010 RTI in Go with Go, Python, and C++
federate paths. Its service implementation is coupled to deterministic event
evidence, generation-aware lifecycle controls, and executable conformance
scenarios. The fair-comparison harness requires complete semantic parity and
delivery accounting before producing performance statistics, making both
positive and negative optimization results auditable.

The software is suitable for education, CI, RTI algorithm research, and
single-node distributed simulation where inspectability and repeatability are
important. The current release does not claim formal certification, Java or
legacy HLA APIs, or production clustering. The source, tests, FOM modules,
schemas, documentation, and verification scripts are released under MIT. An
immutable v0.9 tag and evidence archive DOI will be attached before submission.

\section*{CRediT authorship contribution statement}

The author must confirm the final CRediT roles before submission. Candidate
roles are Conceptualization, Methodology, Software, Validation, Formal
analysis, Investigation, Data curation, Visualization, Writing -- original
draft, and Writing -- review and editing.

\section*{Funding}

Funding information must be confirmed before submission. If no specific
funding supported the work, the final manuscript will use the journal's
recommended no-specific-funding statement.

\section*{Declaration of competing interest}

The author must complete Elsevier's declaration tool and insert the resulting
statement before submission.

\section*{Declaration of generative AI and AI-assisted technologies in the
manuscript preparation process}

During preparation of this draft, the author used OpenAI Codex to help
organize, edit, and validate the repository documentation and manuscript
structure. The author must review and edit the content, verify all evidence and
references, and assumes full responsibility for the submitted work.

\section*{Data availability}

Executable inputs are available in the source repository. The immutable v0.9
source archive and the permissible claim-grade evidence archive, including
manifests, raw samples, checksums, and \texttt{analysis.json}, will be linked
before submission. Proprietary Pitch binaries and restricted logs will not be
redistributed.

\section*{Acknowledgements}

Acknowledgements, if any, must be confirmed before submission.

\begin{thebibliography}{00}

\bibitem{ieee1516}
IEEE Standards Association, IEEE Standard for Modeling and Simulation (M\&S)
High Level Architecture (HLA): Framework and Rules, IEEE Std 1516-2010, 2010.

\bibitem{ieee15161}
IEEE Standards Association, IEEE Standard for Modeling and Simulation (M\&S)
High Level Architecture (HLA): Federate Interface Specification, IEEE Std
1516.1-2010, 2010.

\bibitem{ieee15162}
IEEE Standards Association, IEEE Standard for Modeling and Simulation (M\&S)
High Level Architecture (HLA): Object Model Template, IEEE Std 1516.2-2010,
2010.

\bibitem{fujimoto}
R. M. Fujimoto, Parallel and Distributed Simulation Systems, Wiley, 2000.

\bibitem{gorti}
C. Choi, gorti: An open-source IEEE 1516-2010 Run-Time Infrastructure in Go,
version 0.9.0, software archive DOI to be assigned before submission.

\end{thebibliography}

\section*{Current executable software version}

The executable submission version will be the verified v0.9.0 GitHub release.
Release archives and checksums must be live before this manuscript is
submitted.

\end{document}
\endinput
